\scanpage{44}
\mspage{42.}

\[
 \mathrm{ii}''\Longrightarrow\mathrm{i}\Longrightarrow\mathrm{i}'\deleted{\unsure{\relbar}}
\]

tautologiquement, il reste à prouver que i$'$ implique \inserted{i$''$ et i$''$} (ii$'$). On ne peut appliquer directement la prop. 1, car l'hyp. \deleted{\unsure{n'est}} \inserted{(i$'$)} ne porte que sur les objets de $\mathrm{Ind}\,C$ qui proviennent de $C$, et il n'est pas évident à priori qu'elle passe aux objets quelconques.

Mais cela l'est dans le cas d'un hom de ind-objets sur \unsure{les} $(X_i),(Y_i)$ sur \unsure{le|la} $=$ \unsure{ens|cat} d'indices, \deleted{\unsure{qui est un}} \inserted{donné par} \deleted{\unsure{hom} $X_i\unsure{\relbar|\to}\unsure{Z|Y_i}$} un \deleted{\unsure{syst.}} hom ordinaire (\unsure{terme à terme}) de \deleted{\unsure{inductifs}} systèmes inductifs sur $I$, $X_i\to\unsure{Z_i|Y_i}$ : \deleted{\unsure{Si les}} si les $X_i\unsure{\relbar|\to}Y_i$ sont \unsure{dans|des} $W$, l'hyp. (i) implique que $H(Y)\to H(X)$ est iso, et de $=$ pour i$'$ pour

\[
 \underset{\substack{\Vert\\(X_i)}}{X}\longrightarrow
 \underset{\substack{\Vert\\(\psi\varphi(X_i))_i}}{\Psi\Phi(X)} ,
\]

où l'hom. d'adjonction est donné terme à terme par les hom d'adjonction $X_i\longrightarrow\psi\varphi(X_i)$. Donc l'hyp. i$'$ \deleted{de prop. \unsure{2}} ici implique l'hyp. i$'$ de prop 1 pour $\Phi$ et $\Psi$, donc aussi \unsure{(ii$'$) ici} (par prop. 1), et on gagne.

\figure{\begin{tikzpicture}[x=1cm,y=1cm]\draw (0,0.28)--(0.15,0);\end{tikzpicture}}

Revenons à la situation où $C=\mathrm{Cat\,hot}$, $C'=\mathrm{Hot}$, $\varphi$ le foncteur localisation. Si $K\in\mathrm{Ob}\,\mathrm{Cat\,hot}$, alors $\psi\varphi(K)$ s'appelle \uline{l'enveloppe de \inserted{Thomason-}\unsure{Kan}} de $K$ --- si $K$ est une catégorie, son enveloppe $\overline K$ \deleted{et le} et la \unsure{flèche|flèches} $K\unsure{\relbar|\to}\overline K$ ne sont définis que dans la catégorie Cat hot, i.e. $\overline K$ n'est défini qu'à ``homotopie près'',

% archive sheet number: 43
